
%% Complete JRSI manuscript using the uploaded .tex only as a formatting template
\documentclass[openacc]{rsproca_new}

\usepackage[T1]{fontenc}
\usepackage[utf8]{inputenc}
\usepackage{lmodern}
\usepackage{microtype}
\usepackage{graphicx}
\usepackage{booktabs}
\usepackage{longtable}
\usepackage{array}
\usepackage{amsmath,amssymb}
\usepackage{threeparttable}
\usepackage{multicol}
\usepackage{url}
\usepackage{csquotes}

\jname{rsif}
\Journal{J R Soc Interface\ }

\begin{document}

\title{Dynamic compliance in a balanced pelvic ring: sex-specific tuning of support and childbirth-relevant deformation pathways}

\author{Petr Heny\v{s}$^{1}$, Michal Kucha\v{r}$^{1,2}$, Alexander Mor\'{a}vek$^{2}$, Niels Hammer$^{3,4}$}

\address{$^{1}$Institute of New Technologies and Applied Informatics, Faculty of Mechatronics, Informatics and Interdisciplinary Studies, Technical University of Liberec, Liberec, Czech Republic\\
$^{2}$Department of Anatomy, Faculty of Medicine in Hradec Kr\'{a}lov\'{e}, Charles University, Hradec Kr\'{a}lov\'{e}, Czech Republic\\
$^{3}$Division of Macroscopic and Clinical Anatomy, Gottfried Schatz Research Center, Medical University of Graz, Graz, Austria\\
$^{4}$Division of Macroscopic and Clinical Anatomy, Gottfried Schatz Research Center, Medical University of Graz, Graz, Austria; Department of Orthopedic and Trauma Surgery, University of Leipzig, Leipzig, Germany; Division of Biomechatronics, Fraunhofer IWU, Dresden, Germany}

\subject{biomechanics, evolutionary anthropology, obstetrics}

\keywords{biomechanics, human pelvis, modal analysis, finite-element modelling, childbirth, allometry, sacroiliac joint}

\corres{Michal Kucha\v{r}\\
\email{kucharm@lfhk.cuni.cz}}

\begin{abstract}
The human pelvis must support habitual bipedal loading while also accommodating the rare mechanical demands of childbirth. This dual role is still often framed through static canal dimensions, although the pelvis functions as a compliant ring. We analysed 278 anatomically standardised lumbopelvic finite-element models derived from clinical CT using correspondence-preserving modal decomposition, allometric normalisation, four standardised loading scenarios, and sensitivity analyses of ligamentous and cartilaginous constraints. Pelvic size explained much of the variation in raw eigenvalues, but the sex signal persisted after adjustment for scale and age, most clearly in modes 2 and 9. Under bilateral and unilateral stance, both sexes recruited the same low-order support repertoire, indicating a shared backbone for routine load transfer. Under simulated inlet opening, mode 9 remained dominant in both sexes, but female response was less concentrated in that single mode and more broadly distributed across coordinated low-order pathways. Sensitivity analysis localised this inlet-opening pathway chiefly to the pubic symphysis, anterior sacroiliac ligaments, and sacroiliac cartilage interfaces. The female pelvis is therefore best interpreted not as globally lax, but as a balanced compliant ring that preserves everyday support while selectively tuning childbirth-relevant deformation.
\end{abstract}

\maketitle

\begin{multicols}{2}

\section{Introduction}

The human pelvis must satisfy two demanding functions at once. It transmits loads between the trunk and lower limbs during habitual stance and locomotion, yet it also defines the bony boundary of the birth canal. This dual role has long encouraged a predominantly geometric interpretation of pelvic biology, in which efficient bipedalism is treated as a constraining force and parturition as a widening force \cite{rosenberg2002,trevathan2015,grunstra2023}. In that framing, pelvic diameters become the main currency of function.

That framing remains useful, but it is incomplete. The pelvis is not simply a set of obstetric distances. It is a closed load-bearing ring whose behaviour depends on geometry, material distribution, cartilage interfaces and ligamentous constraints. The mechanically relevant question is therefore not only whether a pelvis is wide or narrow, but also which deformation pathways are available, how accessible they are, and how a given load recruits them.

Several lines of work make this broader view necessary. Wider female pelves do not inevitably impose a locomotor penalty \cite{warrener2015}. Recent evolutionary and obstetric syntheses have argued that childbirth should not be reduced to a one-dimensional opposition between locomotion and canal size \cite{grunstra2023,pavlicev2020}. Human pelvic form also shows substantial geographic and developmental variation \cite{betti2018,verbruggen2017,novak2020}, which makes it unlikely that any single linear dimension can stand in for pelvic function in general.

What remains limited is a population-scale mechanical framework. Subject-specific finite-element studies have shown that sacroiliac constraints, pubic symphysis stiffness and load direction materially affect pelvic behaviour \cite{eichenseer2011,hammer2013}. This is biomechanically plausible because the sacroiliac joint is a specialised load-transfer articulation with limited but functionally meaningful motion \cite{vleeming2012}. Yet isolated case studies do not readily show which features of behaviour are shared across individuals, which vary with scale, and which remain sexually structured after size is taken into account.

Here we address that gap using correspondence-preserving modal analysis of 278 anatomically standardised human pelves. Modal decomposition treats the pelvis not as a rigid frame, but as a repertoire of deformation pathways with different stiffness costs. This allows two related but distinct questions to be separated. First, which pathways are mechanically accessible? Second, once those pathways are available, how is response organised under a given load? That distinction is central to childbirth-related mechanics.

Our working hypothesis is that the female pelvis is best understood not as globally loosened, but as a balanced compliant ring. In this framework, allometry should shape the baseline accessibility landscape, whereas sex should selectively tune how stiffness is distributed across the ring and how low-order pathways are recruited under different loading conditions. We therefore asked four linked questions. How much of the raw spectrum is explained by size? Which sex differences remain after adjustment for scale and age? Do habitual support loads and childbirth-related loads draw on the same modal repertoire? And which interfaces and ligaments regulate the principal inlet-opening pathway?

\section{Material and methods}

\subsection{Cohort and correspondence-preserving framework}

The analysis used the BoneDat resource, a standardised database of 278 clinical lumbopelvic CT scans from a Central European cohort aged 16--91 years \cite{henys2025}. The present workbooks comprised 150 female and 128 male individuals. BoneDat provides segmented anatomy, non-linear registrations to a common template and subject-specific deformation fields, allowing each pelvis to be represented on a shared finite-element domain while preserving subject-specific morphology. This correspondence-preserving design makes direct mode-wise and region-wise comparison possible across the full cohort.

The source scans were acquired without a calibration phantom. Mapped material properties are therefore interpreted here as cohort-wise relative descriptors rather than as direct clinical densitometry. Following the released pipeline, Hounsfield units were converted by phantom-less internal calibration to hydroxyapatite-equivalent density and then mapped to elastic properties \cite{henys2025}.

\subsection{Allometry and static morphology}

Allometric descriptors were taken from \texttt{allometry.xlsx}. The principal size variable was the cube root of total pelvic volume relative to the template reference volume. Additional descriptors included total surface area, total mass, effective thickness and effective density. Static morphological variables were taken from \texttt{anatomy\_data.xlsx} and included the anteroposterior inlet diameter (AP), pelvic inlet transverse diameter (PIT), subpubic angle, biacetabular width, biischiadic width, bituberous width, sacral width and iliopectineal eminence width.

\subsection{Finite-element modelling and modal decomposition}

Each released deformation field was used to warp a tetrahedral template mesh into subject space, after which mechanics were evaluated on the shared reference domain. Bone, sacroiliac cartilage and the pubic symphysis were represented on this common mesh. Bone properties were spatially heterogeneous; cartilage and the pubic symphysis were assigned region-wise properties. Linear isotropic elasticity was used throughout. For each subject, pelvic mechanics were summarised through the generalised eigenproblem $K\phi_k = \lambda_k M\phi_k$, where $K$ and $M$ are the subject-specific stiffness and mass matrices. Lower eigenvalues were interpreted as mechanically more accessible deformation pathways. The present analysis focused on modes 1--10, the biologically relevant low-order spectrum documented in the workbooks.

\subsection{Loading scenarios and modal participation}

Static responses were evaluated for four standardised scenarios available in the released workbooks: bilateral stance (SP2leg), unilateral stance (SP1leg), inlet opening (LAB phase1) and outlet opening (LAB phase2). These scenarios were treated as controlled mechanical probes rather than literal simulations of labour. Static displacement fields were projected onto the modal basis, and modal strain-energy fractions were normalised over modes 1--10. In this article, eigenvalues are interpreted as pathway accessibility, whereas load-case fractions are interpreted as pathway recruitment under a given load.

\subsection{Sensitivity analysis}

To identify which structures regulate the most relevant pathways, we analysed workbook-based sensitivities of eigenvalues to changes in ligament stiffness and interface modulus using \texttt{ligament\_sensitivity.xlsx} and \texttt{cartilage\_sensitivity.xlsx}. The main mechanistic focus was mode 9, because it combined strong sex structure with dominant inlet-opening recruitment. Modes 2 and 5 were retained as support-related and outlet-related reference modes. Pretension perturbations were reviewed separately but were not integrated into the main hierarchy because they probe prestress rather than constitutive stiffness of the base eigenproblem.

\subsection{Statistical analysis}

Static sex differences were assessed with Welch two-sample $t$-tests. Mode-wise sex effects were estimated using linear models of the form $\log(\lambda) \sim \log(\mathrm{scale}) + \mathrm{sex} + \mathrm{age}$, with the sex coefficient expressed as the percentage difference for males relative to females after adjustment for scale and age. Female-only univariate models related size-scaled eigenvalues of modes 9 and 2 to individual pelvic dimensions in order to test how much of the mechanical phenotype could be recovered from conventional pelvimetry.

\section{Results}

\subsection{Cohort structure and static dimorphism}

The cohort reproduced the expected static sex pattern (table~\ref{tab:cohort}). Female pelves were smaller in overall scale but showed larger AP and PIT measures, a wider subpubic angle, and larger biischiadic and bituberous widths. Sacral width did not differ significantly between the sexes.

\subsection{Allometry sets the baseline accessibility landscape}

Allometric regression showed near-isometric scaling of total volume with pelvic scale (slope 3.000, $R^2 = 1.000$), shallower scaling of total surface area (1.684, $R^2 = 0.894$), increasing total mass with scale (2.679, $R^2 = 0.902$), and a modest decline in effective density ($-0.321$, $R^2 = 0.117$; table~\ref{tab:allometry}). In practical terms, this means that raw eigenvalues cannot be interpreted without reference to size. In the present modal framework, smaller geometrically similar structures are expected to show higher raw eigenvalues. That expectation was met across much of the spectrum, but not completely: modes 2 and 9 already departed from it in the unscaled data, with female means lower than male means despite smaller average female pelvic size.

\subsection{Scale-corrected eigenstructure reveals selective sex effects}

After adjustment for scale and age, the strongest sex effects were retained in modes 2 and 9, where male pelves were 20.7\% and 20.1\% stiffer, respectively (figure~\ref{fig:sexeffects}; table~\ref{tab:modes}). Substantial positive male--female differences were also present in modes 4 (15.9\%) and 8 (14.5\%), with smaller but significant effects in modes 3, 7 and 10. By contrast, mode 1 was effectively neutral, and modes 5 and 6 did not retain robust sex effects after adjustment. The residual sex signal was therefore concentrated in selected pathways rather than diffused across the full spectrum.

\subsection{Support and childbirth-related loads recruit a shared repertoire with sex-biased weighting}

Under bilateral stance, mode 2 dominated the response in every individual and accounted for about 85\% of mean modal participation in both sexes (table~\ref{tab:loadcases}; figure~\ref{fig:recruitment}). Unilateral stance broadened recruitment but remained built on the same low-order repertoire, dominated by modes 1 and 2 with a smaller contribution from mode 3. These scenarios therefore point to a largely shared support backbone.

The obstetric probes recruited the same general repertoire less uniformly. Under simulated inlet opening, mode 9 was the dominant component in 265 of 278 individuals. However, the sex difference in this scenario was not a reversal of the dominant pathway. Both sexes remained centred on the same inlet-opening axis, but the weighting differed: women placed less of the mean response into mode 9 itself and more into neighbouring low-order pathways, especially mode 4, whereas men concentrated more strongly in mode 9 and mode 10. Under outlet opening, mode 5 became the dominant component in both sexes, with sex differences appearing mainly in the secondary weighting of modes 6 and 9 rather than in the identity of the dominant pathway.

\subsection{Static pelvimetry explains only a modest fraction of pathway accessibility}

In female-only models, the strongest single predictor of size-scaled mode 9 stiffness was PIT ($R^2 = 0.135$), followed by AP ($R^2 = 0.092$), iliopectineal eminence width ($R^2 = 0.086$) and sacral width ($R^2 = 0.077$; table~\ref{tab:staticassoc}). For mode 2, PIT again provided the strongest single association ($R^2 = 0.089$), followed by bituberous width ($R^2 = 0.078$). No single static dimension explained more than a modest proportion of the variance in either pathway.

\subsection{Mode 9 localises inlet-opening mechanics to anterior ring constraints}

Sensitivity analysis identified mode 9 as the most clearly localised childbirth-relevant pathway (figure~\ref{fig:sensitivity}; table~\ref{tab:sensitivity}). The largest interface sensitivities were found in the left and right sacroiliac cartilage interfaces and the pubic symphysis interface. At the ligament level, the strongest effect was associated with the symphyseal ligament group, followed by the anterior sacroiliac ligaments and then the interosseous ligaments. Posterior sacroiliac, sacrospinous and sacrotuberous ligaments were markedly less influential. Relative to reference modes 2 and 5, mode 9 therefore showed a distinctly anteriorly weighted sensitivity profile.

\section{Discussion}

The central implication of this study is that the apparent conflict between pelvic support and childbirth need not be resolved by exchanging one function for the other. Instead, the results suggest a more nuanced mechanical solution: the pelvis preserves a shared support architecture for everyday load transfer while redistributing the accessibility and load-dependent recruitment of selected deformation pathways under childbirth-related loading. Read in that way, the key result is not static dimorphism alone, but selective tuning within a common compliant ring.

That interpretation depends first on separating size from mechanics. Raw eigenvalues were strongly shaped by scale, as expected for structures that are mechanically similar but differ in overall size. The allometric results therefore matter substantively, not just statistically. Near-isometric scaling of pelvic volume and the strong scale dependence of much of the raw spectrum mean that any claim of greater female compliance based only on unadjusted eigenvalues would risk confounding organisation with size. Once scale and age were taken into account, however, sex effects persisted in selected modes, especially modes 2 and 9. The relevant difference is therefore not simply that female pelves are smaller, but that stiffness is redistributed differently across the same ring system.

The support scenarios show why that distinction is important. Bilateral and unilateral stance drew on the same low-order repertoire in both sexes, with bilateral stance dominated by mode 2 and unilateral stance organised primarily around modes 1--3. This shared backbone argues against interpreting the female pelvis as globally lax or mechanically compromised for routine support. At the same time, the adjusted sex effect in mode 2 indicates that shared does not mean identical. The support architecture is conserved, but it remains selectively tuned.

Conventional pelvic dimensions remain relevant in this framework, but their explanatory reach is limited. The cohort reproduced the expected sex differences in inlet dimensions, subpubic angle and outlet-related breadths, confirming that the sample behaves morphologically as anticipated. Yet those static measures recovered only a modest fraction of the variance in the mechanically important pathways. Even for mode 9, the strongest single association was with PIT, and it explained only a small proportion of the variance. Static dimensions therefore describe part of the structural setting, but they do not by themselves capture how the ring deforms under load. This is precisely where a pathway-based analysis adds value beyond classical pelvimetry.

That added value becomes clearest under childbirth-related loading. Inlet opening remained centred on mode 9 in both sexes, so the female pattern is not explained by switching to a different dominant mechanism. Rather, the sexes differed in how that shared mechanism was weighted. Men concentrated more of the inlet-opening response into mode 9 itself, whereas women distributed more of the response across neighbouring low-order pathways, especially mode 4, while still retaining mode 9 as the principal axis. Outlet opening showed a comparable logic at a different locus: both sexes remained centred on mode 5, with sex differences appearing mainly in the secondary weighting of adjacent modes rather than in the identity of the dominant pathway. These results are more consistent with selective redistribution than with wholesale remodelling of pelvic mechanics.

Mode 9 is therefore central not because women use more of it in fractional terms, but because it sits at the intersection of accessibility and recruitment. It already departs from the raw allometric expectation, it retains one of the strongest adjusted sex effects, and it remains the dominant inlet-opening component in nearly the entire cohort. This combination makes mode 9 the clearest childbirth-relevant deformation pathway in the present framework. The key female feature is not maximal reliance on that one pathway, but broader access to a coordinated neighbourhood around it.

The sensitivity results further constrain that interpretation. Mode 9 was governed chiefly by the pubic symphysis, the sacroiliac cartilage interfaces and the anterior sacroiliac ligaments, whereas posterior sacroiliac, sacrospinous and sacrotuberous ligaments were much less influential. This pattern localises childbirth-relevant tuning to specific ring interfaces and anterior constraint structures rather than to diffuse relaxation of the pelvis as a whole. It also strengthens the argument that the mechanically relevant phenotype is a property of jointed ring organisation, not a simple consequence of canal width. Because all subjects were compared on a correspondence-preserving template, the repeated localisation of mode 9 to the same interfaces is most parsimoniously interpreted as a shared mechanical pathway whose accessibility is differentially tuned across individuals.

Taken together, these findings support a more precise reformulation of the locomotion--childbirth problem. The data do not indicate two fundamentally different pelvic systems in women and men, nor do they suggest that childbirth-relevant compliance is achieved by sacrificing ordinary support function. Instead, the same balanced ring appears to maintain a common support backbone while redistributing the accessibility and recruitment of selected low-order deformation pathways under childbirth-related loading. This interpretation aligns with broader reassessments of the obstetrical dilemma, which argue that childbirth cannot be reduced to canal dimensions alone and that locomotor competence is not necessarily undermined by a wider or differently organised female pelvis \cite{grunstra2023,warrener2015,pavlicev2020,betti2018}. The present study extends that debate by offering a population-scale mechanical formulation of how both demands may be accommodated in the same structure.

Several limitations should remain explicit. The models are linear and confined to small-deformation mechanics, the cohort is retrospective and non-pregnant, and the sample derives from a Central European clinical population rather than an obstetric outcomes cohort. Material mapping relied on phantom-less calibration and is best interpreted comparatively rather than as direct densitometry. The loading scenarios were standardised probes rather than patient-specific labour trajectories, and pretension analyses were kept interpretatively secondary because they address prestress rather than constitutive stiffness in the base eigenproblem. These constraints do not negate the main finding, but they set its present scope.

The next step is to test whether the same pathway structure persists in pregnancy-aware and non-linear models, whether it relates to clinically meaningful obstetric outcomes, and whether similar tuning is seen across more diverse populations. More broadly, the framework developed here suggests that future work should move beyond asking whether a pelvis is simply wide or narrow and instead ask how the pelvic ring makes selected deformation pathways available under different mechanical demands.

\section{Conclusion}

This study shows that the human pelvis is best interpreted as a load-dependent compliant ring. Allometry explains much of the baseline accessibility landscape, but it does not remove the sex signal. After adjustment for scale and age, the strongest residual differences were concentrated in selected low-order modes, especially modes 2 and 9. Habitual support loads recruited a shared backbone in both sexes, whereas childbirth-related loading revealed sex-biased weighting within a shared low-order repertoire. Mode 9 emerged as the clearest inlet-opening pathway and was regulated primarily by the pubic symphysis, anterior sacroiliac ligaments and sacroiliac cartilage interfaces. The female pelvis is therefore not best understood as globally lax, but as a balanced compliant ring that preserves everyday support while selectively tuning childbirth-relevant deformation.

\section*{Ethics}

The analysis used retrospectively acquired clinical CT data released within the BoneDat resource \cite{henys2025}. The final ethics wording should be checked against the dataset paper and local approvals before submission.

\section*{Data accessibility}

The source morphology dataset is available through BoneDat \cite{henys2025}. Derived tables, figure-generation code and manuscript source files should be archived in a public repository before submission; the repository identifier can be inserted at final revision.

\section*{Authors' contributions}

P.H. and M.K. designed the study. P.H. and M.K. performed the computational work. P.H., M.K. and A.M. analysed the data. P.H., M.K. and A.M. drafted the manuscript. N.H. contributed biomechanical and anatomical interpretation and critically revised the text. All authors approved the final version.

\section*{Competing interests}

The authors declare no competing interests.

\section*{Funding}

Funding details should be inserted by the authors before submission.

\section*{Acknowledgements}

We thank the BoneDat project for providing the standardised anatomical framework on which this analysis is based.

\end{multicols}

\begin{table}[p]
\centering
\caption{Cohort summary and static pelvic measures.}
\label{tab:cohort}
\begin{tabular}{lccc}
\toprule
Measure & Female ($n=150$) & Male ($n=128$) & $p$ \\
\midrule
Age (years) & $53.360 \pm 19.736$ & $54.523 \pm 18.912$ & 0.617 \\
Scale & $0.925 \pm 0.045$ & $0.998 \pm 0.045$ & $6.51 \times 10^{-32}$ \\
Effective density & $1.341 \pm 0.079$ & $1.321 \pm 0.074$ & 0.028 \\
AP (mm) & $121.045 \pm 10.017$ & $112.339 \pm 9.070$ & $4.63 \times 10^{-13}$ \\
PIT (mm) & $134.852 \pm 8.590$ & $127.749 \pm 7.605$ & $2.88 \times 10^{-12}$ \\
Subpubic angle (deg) & $85.802 \pm 6.141$ & $69.113 \pm 5.263$ & $7.30 \times 10^{-71}$ \\
Biacetabular width (mm) & $129.076 \pm 8.214$ & $121.976 \pm 7.216$ & $2.94 \times 10^{-13}$ \\
Biischiadic width (mm) & $109.868 \pm 7.128$ & $97.707 \pm 7.057$ & $9.40 \times 10^{-35}$ \\
Bituberous width (mm) & $136.458 \pm 8.291$ & $127.295 \pm 7.492$ & $3.03 \times 10^{-19}$ \\
Sacral width (mm) & $116.746 \pm 5.668$ & $117.473 \pm 6.538$ & 0.327 \\
Iliopectineal eminence width (mm) & $106.203 \pm 7.020$ & $97.902 \pm 5.963$ & $1.90 \times 10^{-22}$ \\
\bottomrule
\end{tabular}
\end{table}

\begin{table}[p]
\centering
\caption{Allometric scaling of geometric and material-derived quantities.}
\label{tab:allometry}
\begin{tabular}{lccc}
\toprule
Outcome versus scale & Slope & $p$ & $R^2$ \\
\midrule
Total volume & 3.000 & $< 1 \times 10^{-300}$ & 1.000 \\
Total surface area & 1.684 & $2.53 \times 10^{-136}$ & 0.894 \\
Total mass & 2.679 & $2.58 \times 10^{-141}$ & 0.902 \\
Effective density & $-0.321$ & $5.07 \times 10^{-9}$ & 0.117 \\
\bottomrule
\end{tabular}
\end{table}

\begin{table}[p]
\centering
\caption{Mode-wise scale dependence and adjusted sex effects (modes 1--10).}
\label{tab:modes}
\begin{tabular}{cccccc}
\toprule
Mode & $r$ (unscaled) & $p$ & $r$ (scaled) & $p$ & Sex effect \\
\midrule
1 & $-0.461$ & $4.96 \times 10^{-16}$ & 0.009 & 0.876 & $0.85\%$ ($p = 0.734$) \\
2 & $-0.331$ & $1.52 \times 10^{-8}$ & 0.374 & $1.20 \times 10^{-10}$ & $20.72\%$ ($p = 2.62 \times 10^{-18}$) \\
3 & $-0.467$ & $1.82 \times 10^{-16}$ & 0.233 & $8.84 \times 10^{-5}$ & $11.23\%$ ($p = 1.80 \times 10^{-7}$) \\
4 & $-0.509$ & $1.04 \times 10^{-19}$ & 0.273 & $3.70 \times 10^{-6}$ & $15.90\%$ ($p = 4.81 \times 10^{-9}$) \\
5 & $-0.607$ & $2.16 \times 10^{-29}$ & $-0.043$ & 0.480 & $-4.18\%$ ($p = 0.106$) \\
6 & $-0.684$ & $9.25 \times 10^{-40}$ & 0.079 & 0.188 & $3.39\%$ ($p = 0.097$) \\
7 & $-0.778$ & $1.44 \times 10^{-57}$ & 0.179 & 0.003 & $5.13\%$ ($p = 9.70 \times 10^{-4}$) \\
8 & $-0.579$ & $2.59 \times 10^{-26}$ & 0.304 & $2.44 \times 10^{-7}$ & $14.45\%$ ($p = 4.88 \times 10^{-11}$) \\
9 & $-0.369$ & $2.06 \times 10^{-10}$ & 0.388 & $1.97 \times 10^{-11}$ & $20.07\%$ ($p = 8.15 \times 10^{-20}$) \\
10 & $-0.630$ & $4.28 \times 10^{-32}$ & 0.218 & $2.56 \times 10^{-4}$ & $8.66\%$ ($p = 1.43 \times 10^{-5}$) \\
\bottomrule
\end{tabular}
\end{table}

\begin{table}[p]
\centering
\caption{Mean modal participation fractions (\%) by loading scenario after normalisation over modes 1--10.}
\label{tab:loadcases}
\begin{tabular}{llrrrrrrrrrr}
\toprule
Scenario & Sex & M1 & M2 & M3 & M4 & M5 & M6 & M7 & M8 & M9 & M10 \\
\midrule
Bilateral stance & Female & 0.4 & 84.6 & 0.1 & 7.3 & 4.0 & 1.0 & 0.1 & 1.3 & 1.1 & 0.1 \\
Bilateral stance & Male & 0.6 & 84.4 & 0.1 & 4.9 & 4.2 & 1.6 & 0.2 & 2.9 & 0.8 & 0.1 \\
Unilateral stance & Female & 39.5 & 37.8 & 15.3 & 2.6 & 1.5 & 1.2 & 0.2 & 0.5 & 0.7 & 0.6 \\
Unilateral stance & Male & 40.3 & 37.1 & 14.1 & 1.9 & 1.4 & 2.3 & 0.2 & 1.0 & 0.7 & 0.9 \\
Inlet opening & Female & 0.5 & 4.1 & 4.2 & 16.9 & 3.5 & 1.5 & 0.9 & 14.3 & 43.0 & 4.0 \\
Inlet opening & Male & 0.7 & 2.8 & 4.7 & 5.4 & 3.1 & 1.4 & 1.0 & 16.5 & 47.7 & 7.7 \\
Outlet opening & Female & 0.3 & 2.6 & 0.4 & 15.6 & 40.3 & 17.7 & 0.9 & 3.8 & 16.0 & 1.6 \\
Outlet opening & Male & 0.4 & 2.3 & 1.0 & 16.1 & 38.5 & 21.7 & 1.2 & 2.7 & 12.2 & 2.1 \\
\bottomrule
\end{tabular}
\end{table}

\begin{table}[p]
\centering
\caption{Female-only associations between static pelvic dimensions and size-scaled eigenvalues.}
\label{tab:staticassoc}
\begin{tabular}{lcccc}
\toprule
& \multicolumn{2}{c}{Mode 9} & \multicolumn{2}{c}{Mode 2} \\
\cmidrule(lr){2-3}\cmidrule(lr){4-5}
Dimension & $R^2$ & $p$ & $R^2$ & $p$ \\
\midrule
PIT & 0.135 & $3.63 \times 10^{-6}$ & 0.089 & $2.11 \times 10^{-4}$ \\
AP & 0.092 & $1.59 \times 10^{-4}$ & 0.002 & 0.600 \\
Iliopectineal eminence width & 0.086 & $2.77 \times 10^{-4}$ & 0.026 & 0.050 \\
Sacral width & 0.077 & $5.96 \times 10^{-4}$ & 0.020 & 0.081 \\
Biacetabular width & 0.012 & 0.173 & 0.053 & 0.005 \\
Biischiadic width & 0.025 & 0.055 & 0.060 & 0.003 \\
Bituberous width & 0.016 & 0.121 & 0.078 & $5.57 \times 10^{-4}$ \\
Subpubic angle & 0.071 & $9.84 \times 10^{-4}$ & 0.000 & 0.845 \\
\bottomrule
\end{tabular}
\end{table}

\begin{table}[p]
\centering
\caption{Sensitivity hierarchy for mode 9, with reference values for modes 5 and 2.}
\label{tab:sensitivity}
\begin{tabular}{lccc}
\toprule
Structure & Mode 9 sensitivity & Mode 5 sensitivity & Mode 2 sensitivity \\
\midrule
Left SIJ cartilage interface & $3.32 \times 10^{-3}$ & $6.63 \times 10^{-4}$ & $2.24 \times 10^{-4}$ \\
Right SIJ cartilage interface & $2.69 \times 10^{-3}$ & $3.74 \times 10^{-4}$ & $1.97 \times 10^{-4}$ \\
Pubic symphysis interface & $2.07 \times 10^{-3}$ & $6.93 \times 10^{-4}$ & $9.45 \times 10^{-6}$ \\
Symphyseal ligament group & $1.09 \times 10^{-6}$ & $7.86 \times 10^{-8}$ & $2.12 \times 10^{-9}$ \\
Anterior SIJ ligaments (bilateral mean) & $4.64 \times 10^{-7}$ & $9.77 \times 10^{-8}$ & $2.59 \times 10^{-8}$ \\
Interosseous SIJ ligaments (bilateral mean) & $2.40 \times 10^{-7}$ & $4.31 \times 10^{-9}$ & $4.68 \times 10^{-11}$ \\
Posterior SIJ ligaments (bilateral mean) & $1.27 \times 10^{-8}$ & $3.47 \times 10^{-8}$ & $2.87 \times 10^{-9}$ \\
Sacrospinous ligaments (bilateral mean) & $6.57 \times 10^{-9}$ & $3.07 \times 10^{-9}$ & $5.85 \times 10^{-9}$ \\
Sacrotuberous ligaments (bilateral mean) & $2.37 \times 10^{-8}$ & $1.24 \times 10^{-9}$ & $6.49 \times 10^{-10}$ \\
\bottomrule
\end{tabular}
\end{table}

\begin{figure}[p]
\centering
\fbox{\parbox[c][7cm][c]{0.9\linewidth}{\centering Placeholder for Figure~1\\
Scale-adjusted sex effects across low-order modes.}}
\caption{Scale-adjusted sex effects across low-order modes. Points show the adjusted male--female stiffness difference for each mode after control for pelvic scale and age; bars show 95\% confidence intervals. Modes 2 and 9 carry the largest sex effects.}
\label{fig:sexeffects}
\end{figure}

\begin{figure}[p]
\centering
\fbox{\parbox[c][7cm][c]{0.9\linewidth}{\centering Placeholder for Figure~2\\
Load-case recruitment across low-order modes.}}
\caption{Load-case recruitment across low-order modes. Habitual support scenarios are built on a shared low-order backbone in both sexes. Inlet opening remains centred on mode 9 in both sexes but is less concentrated in that single mode in women.}
\label{fig:recruitment}
\end{figure}

\begin{figure}[p]
\centering
\fbox{\parbox[c][7cm][c]{0.9\linewidth}{\centering Placeholder for Figure~3\\
Sensitivity hierarchy for the principal inlet-opening pathway.}}
\caption{Sensitivity hierarchy for the principal inlet-opening pathway. Mode 9 is governed chiefly by the pubic symphysis, anterior sacroiliac ligaments and sacroiliac cartilage interfaces, whereas posterior ligament groups are much less influential.}
\label{fig:sensitivity}
\end{figure}

\clearpage

\begin{thebibliography}{99}

\bibitem{rosenberg2002}
Rosenberg K, Trevathan WR. 2002 Birth, obstetrics and human evolution. \textit{BJOG} \textbf{109}, 1199--1206.

\bibitem{trevathan2015}
Trevathan WR. 2015 Primate pelvic anatomy and implications for birth. \textit{Philos. Trans. R. Soc. B} \textbf{370}, 20140065.

\bibitem{grunstra2023}
Grunstra NDS \textit{et al.} 2023 There is an obstetrical dilemma: misconceptions about the evolution of human childbirth and pelvic form. \textit{Am. J. Biol. Anthropol.} \textbf{181}, 535--544.

\bibitem{warrener2015}
Warrener AG, Lewton KL, Pontzer H, Lieberman DE. 2015 A wider pelvis does not increase locomotor cost in humans, with implications for the evolution of childbirth. \textit{PLoS ONE} \textbf{10}, e0118903.

\bibitem{pavlicev2020}
Pavli\v{c}ev M, Romero R, Mitteroecker P. 2020 Evolution of the human pelvis and obstructed labor: new explanations of an old obstetrical dilemma. \textit{Am. J. Obstet. Gynecol.} \textbf{222}, 3--16.

\bibitem{betti2018}
Betti L, Manica A. 2018 Human variation in the shape of the birth canal is significant and geographically structured. \textit{Proc. R. Soc. B} \textbf{285}, 20181807.

\bibitem{verbruggen2017}
Verbruggen SW, Nowlan NC. 2017 Ontogeny of the human pelvis. \textit{Anat. Rec.} \textbf{300}, 643--652.

\bibitem{novak2020}
Nov\'{a}k J, Br\r{u}\v{z}ek J, Zamrazilov\'{a} H, Va\v{n}kov\'{a} M, Hill M, Sedl\'{a}k P. 2020 The relationship between adolescent obesity and pelvis dimensions in adulthood: a retrospective longitudinal study. \textit{PeerJ} \textbf{8}, e8951.

\bibitem{eichenseer2011}
Eichenseer PH, Sybert DR, Cotton JR. 2011 A finite element analysis of sacroiliac joint ligaments in response to different loading conditions. \textit{Spine} \textbf{36}, E1446--E1452.

\bibitem{hammer2013}
Hammer N, Steinke H, Lingslebe U, Bechmann I, Josten C, Slowik V, B\"{o}hme J. 2013 Ligamentous influence in pelvic load distribution. \textit{Spine J.} \textbf{13}, 1321--1330.

\bibitem{vleeming2012}
Vleeming A, Schuenke MD, Masi AT, Carreiro JE, Danneels F, Willard FH. 2012 The sacroiliac joint: an overview of its anatomy, function and potential clinical implications. \textit{J. Anat.} \textbf{221}, 537--567.

\bibitem{henys2025}
Heny\v{s} P, Kucha\v{r} M. 2025 BoneDat, a database of standardised bone morphology for in silico analyses. \textit{Sci. Data} \textbf{12}, 1043.

\end{thebibliography}

\end{document}
